% secciones/marco_conceptual.tex

\section{Marco Conceptual}

\subsection{Big Data y el procesamiento distribuido}

El crecimiento exponencial de los datos generados por empresas, redes sociales, sensores
y plataformas digitales ha impulsado el desarrollo de tecnologías capaces de almacenar
y procesar grandes volúmenes de información. Este fenómeno es conocido como \textit{Big Data}
y suele caracterizarse mediante las cinco V: volumen, velocidad, variedad, veracidad y valor.
En el contexto empresarial, especialmente en el sector retail, millones de transacciones pueden
generarse diariamente, haciendo inviable el procesamiento mediante sistemas tradicionales.

El procesamiento distribuido consiste en dividir una tarea compleja en múltiples subtareas
que pueden ejecutarse simultáneamente sobre diferentes nodos de un clúster. Esta estrategia
permite reducir los tiempos de procesamiento y mejorar la escalabilidad del sistema.
Frameworks como Apache Hadoop y Apache Spark implementan este paradigma para resolver
problemas relacionados con el análisis masivo de datos.

\subsection{Data Warehouse}

Un Data Warehouse es un sistema diseñado para almacenar información histórica y facilitar
el análisis de datos. A diferencia de las bases de datos transaccionales, un Data Warehouse
está optimizado para consultas analíticas y agregaciones. Generalmente se utilizan modelos
dimensionales compuestos por tablas de hechos y tablas dimensión.

En el contexto del retail, una tabla de ventas puede relacionarse con dimensiones como
clientes, productos, tiendas, promociones y fechas. Para este proyecto se utiliza el
esquema estrella definido por el benchmark \tpcds{}, donde \texttt{store\_sales} es la
tabla de hechos central y \texttt{customer}, \texttt{item}, \texttt{store} y
\texttt{date\_dim} son las dimensiones.

\subsection{Benchmark TPC-DS}

\tpcds{} es un benchmark desarrollado por el \textit{Transaction Processing Performance
Council} para evaluar sistemas de Data Warehouse y plataformas analíticas. Simula un
entorno comercial con múltiples canales de venta, incluyendo tiendas físicas, ventas web
y catálogos.

\tpcds{} proporciona un generador de datos capaz de crear datasets desde algunos gigabytes
hasta varios terabytes, permitiendo evaluar el rendimiento de diferentes tecnologías
analíticas bajo escenarios realistas. Para este proyecto se generaron 10 GB de datos
simulando un entorno retail con millones de transacciones.

\subsection{Apache Hive}

Apache Hive es un sistema de Data Warehouse construido sobre Hadoop que permite consultar
y gestionar grandes datasets almacenados en HDFS mediante un lenguaje similar a SQL
denominado HiveQL. Hive traduce las consultas en trabajos MapReduce o Tez, lo que lo hace
adecuado para consultas analíticas batch sobre grandes volúmenes de datos.

\subsection{Apache Spark}

Apache Spark es un motor de procesamiento distribuido en memoria diseñado para cargas de
trabajo analíticas y de machine learning. Spark SQL permite ejecutar consultas SQL sobre
DataFrames distribuidos, ofreciendo un rendimiento significativamente superior a MapReduce
en la mayoría de escenarios al reducir las operaciones de lectura y escritura en disco.

\subsection{Análisis Agéntico de datos}

El análisis agéntico es un paradigma que combina modelos de lenguaje, herramientas
analíticas y mecanismos de automatización para facilitar la exploración y consulta de
grandes volúmenes de datos. En lugar de escribir consultas manualmente, el usuario puede
expresar sus necesidades en lenguaje natural, mientras el sistema interpreta la intención,
determina las operaciones necesarias y utiliza las herramientas adecuadas para obtener
los resultados.

A diferencia de un chatbot convencional, un sistema de análisis agéntico no se limita a
generar texto, sino que puede interactuar con fuentes de datos, ejecutar consultas y
combinar múltiples procesos para resolver problemas específicos del dominio de análisis.

Para este proyecto se implementa una capa agéntica basada en el paradigma de \textit{skills},
donde cada habilidad implementa una función específica del pipeline:
identificación de intención, generación de SQL mediante la API de Gemini,
ejecución sobre Spark SQL y presentación de resultados.
